% !PRINTED 308
\weberpage{308}
%<<CONT
\weq{2}{
\begin{aligned}
\tau &= \begin{pmatrix} 1, & 0, & 1 \\ 1, & 0, & 0 \\ 0, & 1, & 0 \end{pmatrix}, &
\tau^2 &= \begin{pmatrix} 1, & 1, & 1 \\ 1, & 0, & 1 \\ 1, & 0, & 0 \end{pmatrix}, &
\tau^3 &= \begin{pmatrix} 0, & 1, & 1 \\ 1, & 1, & 1 \\ 1, & 0, & 1 \end{pmatrix}. \\
\tau^4 &= \begin{pmatrix} 1, & 1, & 0 \\ 0, & 1, & 1 \\ 1, & 1, & 1 \end{pmatrix}, &
\tau^5 &= \begin{pmatrix} 0, & 0, & 1 \\ 1, & 1, & 0 \\ 0, & 1, & 1 \end{pmatrix}, &
\tau^6 &= \begin{pmatrix} 0, & 1, & 0 \\ 0, & 0, & 1 \\ 1, & 1, & 0 \end{pmatrix}
\end{aligned}
}

Next one selects, among the elements of degree 2, of which
there are 21, a suitable one for $\omega$; this one, for example, is suitable

\weq{3}{\omega = \begin{pmatrix} 1, & 0, & 0 \\ 0, & 1, & 0 \\ 0, & 1, & 1 \end{pmatrix};}

and then, according to the formulae \S.\,72, (17), the elements $\Theta$
and $\chi$ are determined:

\weq{4}{\Theta = \begin{pmatrix} 1, & 1, & 1 \\ 0, & 1, & 0 \\ 0, & 1, & 1 \end{pmatrix}, \quad \Theta^2 = \begin{pmatrix} 1, & 1, & 0 \\ 0, & 1, & 0 \\ 0, & 0, & 1 \end{pmatrix}, \quad \Theta^3 = \begin{pmatrix} 1, & 0, & 1 \\ 0, & 1, & 0 \\ 0, & 1, & 1 \end{pmatrix}}

\weq{5}{\chi = \begin{pmatrix} 1, & 0, & 0 \\ 0, & 0, & 1 \\ 0, & 1, & 1 \end{pmatrix}, \quad \chi^2 = \begin{pmatrix} 1, & 0, & 0 \\ 0, & 1, & 1 \\ 0, & 1, & 0 \end{pmatrix}.}

It is then a matter of a simple calculation to verify the formulae
of Theorem I., \S.\,72, whereby the agreement
of the groups $S$ and $L_7$ is established.

From this it follows that the group $S$, like $L_7$, is \gesperrt{simple}.

\begin{center}
\S.\,80.

\gesperrt{Resolvents of the Equation of Degree 8.}
\end{center}

In order now to make the application to the theory of equations
of degree 8, we denote as before by $Q$ the
cyclic group of degree 8: % [sic] the displayed translation group is elementary abelian of order 8, not cyclic.

\[
\begin{pmatrix} z_1, & z_2, & z_3 \\ z_1 + h_1, & z_2 + h_2, & z_3 + h_3 \end{pmatrix} \text{ (mod 2)},
\]

or more briefly

\[
(z_i,\ z_i + h_i),
\]

and by $S$ the ternary congruence group considered in \S.\,79,
and we set

\[
R = S\,Q,
\]
%CONT>>
